\DocumentMetadata{
  pdfversion=2.0,
  pdfstandard=ua-2,
  lang=en-US,
  tagging=on,
  tagging-setup={math/setup=mathml-SE,tabsorder=structure}
}
\documentclass[11pt,letterpaper]{article}
\usepackage[margin=0.68in,includefoot]{geometry}
\usepackage{newcomputermodern}
\usepackage{amsmath,amscd,mathtools}
\usepackage{tikz,adjustbox}
\providecommand{\square}{\mdlgwhtsquare}
\usepackage[hidelinks]{hyperref}
\hypersetup{
  pdftitle={Combinatorics PhD Exam, September 1991},
  pdfauthor={University of Florida Department of Mathematics},
  pdfsubject={Graduate mathematics examination},
  pdfdisplaydoctitle=true,
  bookmarksopen=true
}
\setcounter{secnumdepth}{0}
\setlength{\parindent}{0pt}
\setlength{\parskip}{0.28em}
\setlength{\emergencystretch}{3em}
\setlength{\leftmargini}{2.15em}
\setlength{\leftmarginii}{2.65em}
\setlength{\leftmarginiii}{3em}
\pagestyle{empty}
\raggedbottom
\allowdisplaybreaks
\NewTaggingSocketPlug{block/list/label}{examproblem}{%
  \tagstructbegin{tag=Lbl}\tagstructend%
  \tagstructbegin{tag=\LBody}%
  \tagstructbegin{tag=\ProblemHeadingTag,title-o={\ProblemHeadingText}}%
  \tagmcbegin{tag=\ProblemHeadingTag,actualtext={\ProblemHeadingText}}%
  #2%
  \tagmcend%
  \tagstructend%
}
\newenvironment{examproblems}[2]{%
  \def\ProblemHeadingTag{#2}%
  \AssignTaggingSocketPlug{block/list/label}{examproblem}%
  \begin{enumerate}%
  \setcounter{enumi}{#1}%
  \renewcommand{\labelenumi}{\textbf{\theenumi.}}%
  \setlength{\topsep}{0.35em}%
  \setlength{\partopsep}{0pt}%
  \setlength{\itemsep}{0.48em}%
  \setlength{\parsep}{0.18em}%
}{\end{enumerate}}
\newenvironment{examsubparts}{%
  \AssignTaggingSocketPlug{block/list/label}{default}%
  \begin{enumerate}%
  \setlength{\topsep}{0.18em}%
  \setlength{\partopsep}{0pt}%
  \setlength{\itemsep}{0.16em}%
  \setlength{\parsep}{0.1em}%
}{\end{enumerate}}
\newenvironment{examinstructions}{%
  \par\begingroup\itshape
}{%
  \par\endgroup\vspace{0.18em}
}
\makeatletter
\renewcommand\subsection{\@startsection{subsection}{2}{0pt}%
  {-1.25ex plus -.25ex minus -.1ex}%
  {0.55ex plus .1ex}%
  {\normalfont\large\bfseries}}
\renewcommand{\@maketitle}{%
  \newpage\null\vskip -1em%
  {\footnotesize\bfseries
    \href{https://www.ufl.edu/}{University of Florida}, \href{https://math.ufl.edu/}{Department of Mathematics}\par}%
  \vskip -0.5em%
  \tagstructbegin{tag=H1,title={Combinatorics PhD Exam, September 1991}}%
  \tagpdfparaOff%
  \tagmcbegin{tag=H1}%
    {\LARGE\bfseries\href{https://gma.math.ufl.edu/exams/combinatorics-phd/}{Combinatorics PhD Exam}, September 1991\par}%
  \tagmcend%
  \tagpdfparaOn%
  \tagstructend%
  \vskip 0.25em%
  \tagmcbegin{artifact}%
    \hrule height 1pt%
  \tagmcend%
  \vskip 1em%
}
\makeatother
\title{Combinatorics PhD Exam, September 1991}
\author{}
\date{}
\begin{document}
\maketitle
\thispagestyle{empty}
\begin{examinstructions}
Do 4 out of 5 problems.
\end{examinstructions}
\begin{examproblems}{0}{H2}
\def\ProblemHeadingText{Problem 1.}
\item
A~\(k\)-arc of a projective plane~\(\Pi\) is a set of~\(k\) points, no three of which are collinear. Prove that every 4-arc in the plane \(PG(2,4)\) lies in exactly two 5-arcs.
\def\ProblemHeadingText{Problem 2.}
\item
Use (without proof) the assertion of problem \#1 to prove that the 6-arcs of \(\Pi=PG(2,4)\) form a block design~\(\Sigma\) on the point set of~\(\Pi\), with \(v=21\) and \(k=6\). Compute~\(\lambda\).
\def\ProblemHeadingText{Problem 3.}
\item
State the Hall Multiplier Theorem, and use it to find a \((91,10,1)\) difference set.
\def\ProblemHeadingText{Problem 4.}
\item
HC is the problem of deciding whether a graph has a Hamiltonian circuit or not, and this problem is known to be NP-complete. Use this fact to prove that HP is also NP-complete, where HP is the problem of deciding whether a graph has a Hamiltonian path (with unspecified initial and terminal vertices).
\def\ProblemHeadingText{Problem 5.}
\item
Let~\(E\) be a finite set and~\(P\) a partition of~\(E\). Call a subset~\(I\) of~\(E\) independent (\(I\in\mathcal I\)) if no two elements of~\(I\) are in the same block of~\(P\).
\begin{examsubparts}
\item[\textnormal{1.}]
Prove that \((E,\mathcal I)\) is a matroid (called a partition matroid).
\item[\textnormal{2.}]
For a bipartite graph~\(B\) with edge set~\(E\), let \(\mathcal M\subseteq 2^E\) denote the set of matchings on~\(B\). Show that \((E,\mathcal M)\) is not a matroid.
\item[\textnormal{3.}]
Prove that \(\mathcal M\) is the intersection of two partition matroids. In other words there are partition matroids \((E,\mathcal I_1)\) and \((E,\mathcal I_2)\) such that \(\mathcal M=\mathcal I_1\cap\mathcal I_2\).

\par
Part 3 is significant because there is a general result that states that if \((E,\mathcal M)\) is such a matroid intersection then there is a polynomial algorithm to find a maximum cardinality independent set.
\end{examsubparts}
\end{examproblems}
\end{document}
